\documentclass[11pt,oneside]{book}

\usepackage[T1]{fontenc}
\usepackage[utf8]{inputenc}
\usepackage{amsmath,amssymb,mathtools}
\usepackage{booktabs}
\usepackage[margin=1in]{geometry}
\usepackage{xcolor}
\usepackage{listings}
\usepackage[hidelinks]{hyperref}

\definecolor{codeblue}{RGB}{25,75,135}
\definecolor{codegreen}{RGB}{25,115,65}
\definecolor{codegray}{RGB}{95,95,95}
\definecolor{codebackground}{RGB}{247,247,247}

\lstdefinestyle{python}{
  language=Python,
  basicstyle=\ttfamily\small,
  keywordstyle=\color{codeblue}\bfseries,
  commentstyle=\color{codegreen},
  stringstyle=\color{codegray},
  backgroundcolor=\color{codebackground},
  frame=single,
  rulecolor=\color{black!20},
  numbers=left,
  numberstyle=\tiny\color{codegray},
  numbersep=8pt,
  showstringspaces=false,
  breaklines=true,
  breakatwhitespace=false,
  columns=fullflexible,
  keepspaces=true,
  tabsize=4
}
\lstset{style=python}

\newcommand{\Hubbard}{\mathcal{H}}
\newcommand{\basisdim}{\mathcal{D}}
\newcommand{\ket}[1]{\left\lvert #1 \right\rangle}
\newcommand{\bra}[1]{\left\langle #1 \right\rvert}
\newcommand{\braket}[2]{\left\langle #1 \middle\vert #2 \right\rangle}
\newcommand{\expect}[1]{\left\langle #1 \right\rangle}
\newcommand{\norm}[1]{\left\lVert #1 \right\rVert}
\newcommand{\code}[1]{\texttt{#1}}

\hypersetup{
  pdftitle={Building a Trustworthy Hubbard Exact-Diagonalization Solver},
  pdfauthor={hubbard-ed contributors},
  pdfsubject={Exact diagonalization tutorial for the one-dimensional Hubbard model}
}


\title{Building a Trustworthy Hubbard\\Exact-Diagonalization Solver\\
\large A step-by-step implementation tutorial}
\author{The \code{hubbard-ed} project}
\date{\today}

\begin{document}

\frontmatter
\hypersetup{pageanchor=false}
\maketitle
\cleardoublepage
\hypersetup{pageanchor=true}
\setcounter{page}{1}

\chapter*{How to use this tutorial}
\addcontentsline{toc}{chapter}{How to use this tutorial}

This tutorial develops a small exact-diagonalization solver for the
one-dimensional spin-$1/2$ Hubbard model.  The order is deliberate: each
chapter adds one abstraction only after the previous layer has a numerical
check.  The result is useful both as an educational implementation and as a
reference against which more ambitious representations can be tested.

Every complete program printed in this document lives in the adjacent
\code{code/} directory and is included with \code{\textbackslash
lstinputlisting}.  The manuscript therefore displays the same source that is
executed during validation.  From the repository root, run

\begin{verbatim}
for script in tutorial/code/step*.py; do
    .venv/bin/python "$script"
done
\end{verbatim}

The tutorial assumes familiarity with elementary second quantization and
linear algebra, but not with a many-body software library.  NumPy and SciPy are
the only numerical dependencies.

\tableofcontents

\mainmatter
\chapter{Model, sectors, and bit-string bases}
\label{chap:basis}

\section{The reference problem}

For a chain of $L$ sites, the spin-$1/2$ Hubbard Hamiltonian is
\begin{equation}
H = -t\sum_{\langle i,j\rangle,\sigma}
\left(c^{\dagger}_{i\sigma}c_{j\sigma}
    +c^{\dagger}_{j\sigma}c_{i\sigma}\right)
+U\sum_{i=0}^{L-1}n_{i\uparrow}n_{i\downarrow}.
\label{eq:hubbard}
\end{equation}
The solver supports open and periodic nearest-neighbor chains.  Both the
up-spin number $N_\uparrow$ and down-spin number $N_\downarrow$ commute with
$H$, so the full Fock space splits into independent blocks.  We solve one block
at a time.

Working in a fixed sector is the first important reduction.  Instead of a
$4^L$-dimensional site Fock space, the sector dimension is
\begin{equation}
\basisdim(L,N_\uparrow,N_\downarrow)
= \binom{L}{N_\uparrow}\binom{L}{N_\downarrow}.
\label{eq:dimension}
\end{equation}
The reduction is exact because the Hamiltonian never changes either particle
number.

\section{Two integers per basis state}

A basis state is represented by
\begin{equation}
s=(b_\uparrow,b_\downarrow),
\end{equation}
where bit $i$ of $b_\sigma$ is one exactly when site $i$ contains a spin
$\sigma$ electron.  Site zero is the least significant bit.  For example, at
$L=4$,
\begin{equation}
(b_\uparrow,b_\downarrow)=(0101_2,0010_2)
\end{equation}
has up electrons on sites $0$ and $2$, and a down electron on site $1$.

This representation makes the common operations direct:
\begin{align}
n_{i\uparrow}(s) &= (b_\uparrow \mathbin{\texttt{>>}} i)\mathbin{\texttt{\&}}1,\\
N_\uparrow(s) &= \operatorname{popcount}(b_\uparrow),\\
N_{\mathrm{double}}(s) &=
\operatorname{popcount}(b_\uparrow\mathbin{\texttt{\&}}b_\downarrow).
\end{align}
Python's integer \code{bit\_count()} supplies the popcount operation.

\section{Enumeration and lookup}

All length-$L$ integers with popcount $N_\uparrow$ are combined with all
length-$L$ integers with popcount $N_\downarrow$.  A tuple stores the resulting
states in deterministic order, while a dictionary maps each state back to its
integer basis index.  The tuple gives $O(1)$ index-to-state lookup and the
dictionary gives average $O(1)$ state-to-index lookup.

The following first program verifies the combinatorial dimension, particle
counts, and round-trip lookup.

\lstinputlisting[caption={Enumerating a fixed sector.},label={lst:basis}]{code/step01_basis.py}

\section{Validation before allocation}

The basis constructor rejects $L<1$, particle numbers outside $[0,L]$, and
noninteger inputs.  It computes Eq.~\eqref{eq:dimension} before allocating the
state tuple.  A configurable dimension guard then refuses unreasonable
sectors.  This ordering matters: a friendly error is better than allowing the
operating system to terminate an overcommitted process.

At half filling with $N_\uparrow=N_\downarrow=L/2$, growth is already severe:
\begin{center}
\begin{tabular}{rrr}
\toprule
$L$ & $(N_\uparrow,N_\downarrow)$ & $\basisdim$ \\
\midrule
6  & $(3,3)$ & 400 \\
8  & $(4,4)$ & 4,900 \\
10 & $(5,5)$ & 63,504 \\
12 & $(6,6)$ & 853,776 \\
\bottomrule
\end{tabular}
\end{center}
The Hamiltonian need not be dense, but the basis and reverse-lookup dictionary
must still fit in memory.

\paragraph{Checkpoint.}
Before writing any operator code, verify for several sectors that every state
has the requested two popcounts, every state is unique, both lookup directions
round-trip, and the dimension equals Eq.~\eqref{eq:dimension}.

\chapter{Fermionic operators and signs}
\label{chap:operators}

\section{One global ordering}

Fermionic signs are meaningless until the spin orbitals have a global order.
This project fixes
\begin{equation}
(0\uparrow,1\uparrow,\ldots,(L-1)\uparrow,
 0\downarrow,1\downarrow,\ldots,(L-1)\downarrow).
\label{eq:ordering}
\end{equation}
Thus site $i$ has global orbital index $i$ for up spin and $L+i$ for down
spin.  The two state integers can be viewed as one $2L$-bit integer,
\begin{equation}
B=b_\uparrow+(b_\downarrow\mathbin{\texttt{<<}}L).
\end{equation}

For orbital $p$, define the parity count
\begin{equation}
m_p(B)=\operatorname{popcount}\left(B\mathbin{\texttt{\&}}(2^p-1)\right).
\end{equation}
It counts occupied orbitals strictly before $p$.  The elementary actions are
\begin{align}
c_p\ket{B} &=( -1)^{m_p(B)}n_p\ket{B\text{ with }p\text{ cleared}},\\
c_p^\dagger\ket{B} &=( -1)^{m_p(B)}(1-n_p)
\ket{B\text{ with }p\text{ set}}.
\end{align}
An invalid annihilation or creation returns no state rather than a zero-amplitude
state.

\section{A hop is two operator actions}

The hopping operator $c_p^\dagger c_q$ acts from right to left:
\begin{enumerate}
\item annihilate at $q$ using the original occupation;
\item create at $p$ using the intermediate occupation;
\item multiply the two signs.
\end{enumerate}
Computing a sign from a shortcut tied only to geometric distance is fragile.
The two-action algorithm remains correct for different spins, long-range
hopping, and periodic boundary bonds.

Consider $L=4$, $b_\uparrow=1010_2$, and the periodic hop $3\to0$.  Site $1$
is occupied between orbitals $0$ and $3$, so
\begin{equation}
c_{0\uparrow}^\dagger c_{3\uparrow}ket{1010_2}=-\ket{0011_2}.
\end{equation}
The sign is not $+1$.  This is the boundary case most likely to expose an
incorrect nearest-neighbor shortcut.

\lstinputlisting[caption={Inspecting creation and periodic hopping signs.},label={lst:signs}]{code/step02_fermion_signs.py}

\section{Algebraic tests}

The implementation is checked against the canonical anticommutators
\begin{align}
\{c_p,c_q^\dagger\}&=\delta_{pq},\\
\{c_p,c_q\}&=0,\\
\{c_p^\dagger,c_q^\dagger\}&=0.
\end{align}
For small $L$, these identities can be tested exhaustively on the complete
$2^{2L}$-state Fock space.  Exhaustive tests are inexpensive here and stronger
than testing only a few hand-selected signs.

The down-spin sign includes all occupied up orbitals because the down block
comes second in Eq.~\eqref{eq:ordering}.  In a number-conserving down-spin hop,
the up-block contribution appears twice and cancels, but creation and
annihilation must still each obey the global convention.  This distinction
becomes important as soon as non-number-conserving operators or Green
functions are introduced.

\paragraph{Checkpoint.}
Do not construct a Hamiltonian until the full small-space anticommutators,
cross-spin signs, an adjacent hop, and an occupied periodic-boundary hop all
pass.

\chapter{Assembling the sparse Hamiltonian}
\label{chap:sparse}

\section{Separate diagonal and hopping actions}

For a basis state $s=(b_\uparrow,b_\downarrow)$, the interaction is diagonal:
\begin{equation}
\bra{s}H_U\ket{s}=U\,
\operatorname{popcount}(b_\uparrow\mathbin{\texttt{\&}}b_\downarrow).
\label{eq:interactiondiag}
\end{equation}
Every allowed hopping term produces a target state in the same
$(N_\uparrow,N_\downarrow)$ sector.  Its matrix element is $-t$ times the sign
returned by the operator layer.

Open chains contain the unique undirected bonds
\begin{equation}
(0,1),(1,2),\ldots,(L-2,L-1).
\end{equation}
Periodic chains add $(L-1,0)$ when $L>2$.  The $L=2$ periodic chain deliberately
contains the same single edge as the open chain; adding it twice would double
the hopping and invalidate the standard two-site analytic result.  For every
undirected edge, both hopping directions and both spin species are applied.

\section{Column-wise construction}

Sparse assembly follows the definition of a matrix action.  For every basis
column $j$ representing $\ket{s_j}$:
\begin{enumerate}
\item append the interaction value from Eq.~\eqref{eq:interactiondiag} at
      $(j,j)$ if it is nonzero;
\item apply each directed spin-preserving hop;
\item map the resulting state $s_i$ to a row index $i$;
\item append the signed matrix element at $(i,j)$.
\end{enumerate}
The triplets $(i,j,H_{ij})$ are first accumulated as a coordinate matrix and
then converted to compressed sparse row (CSR) form.  Duplicate entries are
summed and explicit zeros are removed.

\lstinputlisting[caption={Building and checking a CSR Hamiltonian.},label={lst:sparse}]{code/step03_sparse_hamiltonian.py}

\section{Hermiticity as an implementation invariant}

For real $t$ and $U$, the result must be real symmetric.  More generally,
\begin{equation}
H=H^\dagger.
\end{equation}
This test catches missing reverse hops, inconsistent signs, incorrectly
duplicated bonds, and many indexing errors.  Check both open and periodic
boundaries and more than one particle sector.

Sparse storage changes the scaling from $O(\basisdim^2)$ dense entries to
roughly $O(zL\basisdim)$ possible transitions, where $z$ is a small
coordination-dependent factor.  However, Python lists used during coordinate
assembly have nontrivial overhead.  The implementation therefore has a
separate conservative dimension guard for explicit CSR construction.

\paragraph{Checkpoint.}
For tiny systems, compare the CSR matrix with its conjugate transpose exactly
up to numerical roundoff, inspect the atomic-limit diagonal, and verify at
least one periodic-boundary matrix element whose fermionic sign is negative.

\chapter{Matrix-free action and Lanczos solving}
\label{chap:lanczos}

\section{Apply the same terms without storing them}

Iterative eigensolvers only require products $y=Hx$.  For a vector
$x_j=\braket{s_j}{\psi}$, the matrix-free loop is the sparse assembly loop with
one change: instead of appending $H_{ij}$, accumulate
\begin{equation}
y_i \mathrel{+}= H_{ij}x_j.
\end{equation}
The basis and state-index dictionary are retained, but the CSR arrays are not.
Correctness is checked by applying both implementations to the same random
real and complex vectors.

\lstinputlisting[caption={CSR versus matrix-free multiplication.},label={lst:matrixfree}]{code/step04_matrix_free.py}

Using the same elementary hopping routine in both paths is intentional.  It
reduces the risk that an optimized implementation and a reference
implementation silently adopt different sign conventions.

\section{Lanczos through a linear operator}

SciPy's \code{LinearOperator} wraps the matrix-free action and supplies it to
\code{eigsh}.  For a Hermitian Hamiltonian, requesting
\code{which="SA"} returns the algebraically smallest eigenvalues.  A
deterministic random initial vector is preferable to a constant vector because
a constant vector may lie entirely within one symmetry block and miss a lower
state in another block.

ARPACK requires the number of requested eigenvalues $k$ to satisfy
$k<\basisdim$.  Tiny sectors and complete-spectrum requests below a strict
size threshold therefore use dense diagonalization.  Dense fallback is an
edge-case convenience, not a route for production sectors.

\section{Residuals are the final contract}

An eigensolver's convergence flag is not enough.  Normalize every returned
vector and recompute
\begin{equation}
r_a=\norm{H\psi_a-E_a\psi_a}_2
\label{eq:residual}
\end{equation}
with the matrix-free action.  The default ARPACK tolerance is $10^{-10}$;
typical absolute residuals in the tutorial examples are between $10^{-15}$
and $10^{-10}$.

\lstinputlisting[caption={Low-energy states and independent residual checks.},label={lst:groundstate}]{code/step05_ground_state.py}

Degenerate eigenvalues require care.  Individual eigenvectors within a
degenerate subspace are not unique, so compare the eigenvalue multiset or the
projector onto the subspace rather than component-by-component vectors.

\paragraph{Checkpoint.}
Test CSR and matrix-free products on seeded random real and complex vectors,
compare their low eigenvalues, and require a small independently recomputed
residual for every eigenpair used in a numerical assertion.

\chapter{Observables from the ground state}
\label{chap:observables}

Let
\begin{equation}
\ket{\psi}=\sum_s \psi_s\ket{s},\qquad
\sum_s|\psi_s|^2=1.
\end{equation}
The observable routines normalize any finite nonzero input vector before
measurement.  This is convenient, but solver output should already have unit
norm; a large normalization correction would indicate an upstream problem.

\section{Diagonal observables}

Any occupation-only observable $O(s)$ is diagonal in the site basis:
\begin{equation}
\expect{O}=\sum_s |\psi_s|^2 O(s).
\end{equation}
This covers the total and local double occupancy,
\begin{align}
D &= \sum_i\expect{n_{i\uparrow}n_{i\downarrow}},\\
D/L &= \text{double occupancy per site},
\end{align}
as well as local charge and magnetization,
\begin{align}
\expect{n_i} &= \expect{n_{i\uparrow}+n_{i\downarrow}},\\
\expect{m_i} &= \expect{n_{i\uparrow}-n_{i\downarrow}}.
\end{align}
The project reports $m_i=2S_i^z$ for local magnetization and uses
\begin{equation}
S_i^z=\frac{n_{i\uparrow}-n_{i\downarrow}}{2}
\end{equation}
inside spin correlations.

The longitudinal spin and connected charge correlations are
\begin{align}
C^{zz}_{ij} &= \expect{S_i^zS_j^z},\\
C^n_{ij} &= \expect{n_in_j}-\expect{n_i}\expect{n_j}.
\end{align}
For fixed total particle number, summing $C^n_{ij}$ over both sites gives zero;
this is a useful sum rule.

\section{The one-body density matrix}

The spin-resolved one-body density matrix is not diagonal in the site basis:
\begin{equation}
\gamma^{(\sigma)}_{ij}
=\expect{c_{i\sigma}^\dagger c_{j\sigma}}.
\label{eq:obdm}
\end{equation}
For every source basis state, apply $c_{i\sigma}^\dagger c_{j\sigma}$ using the
same sign routine as the Hamiltonian, find the target basis state, and
accumulate
\begin{equation}
(\psi_{s'})^*\,\mathrm{sign}(s'\leftarrow s)\,\psi_s.
\end{equation}
The basic checks are
\begin{equation}
\gamma^{(\sigma)}=\left(\gamma^{(\sigma)}\right)^\dagger,
\qquad
\operatorname{Tr}\gamma^{(\sigma)}=N_\sigma.
\end{equation}
The eigenvectors of $\gamma$ are natural orbitals, but using them as a basis
for the interacting Hamiltonian requires the transformed interaction discussed
in Chapter~\ref{chap:generalhopping}.

\section{Fourier diagnostics}

On momenta $k_m=2\pi m/L$, the convention for the momentum distribution is
\begin{equation}
n_\sigma(k)=\frac{1}{L}\sum_{ij}e^{ik(i-j)}
\gamma^{(\sigma)}_{ij}.
\end{equation}
It obeys $\sum_k n_\sigma(k)=N_\sigma$.  The charge and longitudinal-spin
structure factors use the same transform:
\begin{align}
N(q)&=\frac{1}{L}\sum_{ij}e^{iq(i-j)}C^n_{ij},\\
S^{zz}(q)&=\frac{1}{L}\sum_{ij}e^{iq(i-j)}C^{zz}_{ij}.
\end{align}
The charge factor is connected by definition.  The spin factor is unconnected
by default and has an explicit connected option.

\lstinputlisting[caption={Measuring site, one-body, and Fourier observables.},label={lst:observables}]{code/step06_observables.py}

\paragraph{Checkpoint.}
Verify particle-number sums, Hermiticity and trace of the one-body density
matrix, reflection symmetry on an open chain, vanishing connected charge
$q=0$ weight, and nonnegative structure factors up to roundoff.

\chapter{Exact results and limiting-case validation}
\label{chap:validation}

A reference solver earns trust by failing loudly when any layer disagrees with
an independent result.  The validation suite is deliberately redundant:
operator algebra, matrix properties, analytic spectra, limiting cases, and
iterative residuals constrain different classes of mistakes.

\section{The validation ladder}

\begin{center}
\begin{tabular}{p{0.24\textwidth}p{0.66\textwidth}}
\toprule
Check & What it detects \\
\midrule
Basis dimension & Missing, duplicated, or out-of-sector states \\
Anticommutators & Global-ordering and parity errors \\
Hermiticity & Missing reverse hops and inconsistent signs \\
Two-site formula & Interaction normalization and hopping scale \\
Full $U=0$ spectrum & Bond convention, degeneracies, and spin factorization \\
Atomic limit & Double-occupancy diagonal \\
Large-$U$ gap & Low-energy scale and excited-state ordering \\
CSR versus matrix-free & Divergent implementations \\
Residual norms & Eigensolver convergence and returned-vector integrity \\
\bottomrule
\end{tabular}
\end{center}

\section{Half-filled two-site model}

For $L=2$, $N_\uparrow=N_\downarrow=1$, and one unique bond, the ground-state
energy is
\begin{equation}
E_0=\frac{U}{2}-\frac{1}{2}\sqrt{U^2+16t^2}.
\label{eq:twosite}
\end{equation}
This single formula tests the relative factors in both Hamiltonian terms.  It
should be checked at several positive and negative $U$ values and more than one
$t$.

At large repulsive $U$, the triplet lies at zero energy while the singlet is
lowered, giving
\begin{equation}
\Delta_{ST}=-E_0
=\frac{\sqrt{U^2+16t^2}-U}{2}
=\frac{4t^2}{U}+O(U^{-3}).
\end{equation}
Thus the solver recovers the superexchange scale $J=4t^2/U$.

\section{Noninteracting limit}

At $U=0$, diagonalize the one-particle hopping problem independently.  For an
open chain,
\begin{equation}
\epsilon_m=-2t\cos\left(\frac{m\pi}{L+1}\right),
\qquad m=1,\ldots,L.
\end{equation}
For a periodic chain with $L\ge3$,
\begin{equation}
\epsilon_m=-2t\cos\left(\frac{2\pi m}{L}\right),
\qquad m=0,\ldots,L-1.
\end{equation}
Under the unique-edge convention, $L=2$ instead has energies $(-|t|,|t|)$.

Every many-body energy in the fixed sector is a sum of occupied one-particle
energies,
\begin{equation}
E=\sum_{a\in A_\uparrow}\epsilon_a
 +\sum_{b\in A_\downarrow}\epsilon_b.
\end{equation}
Comparing the complete small-system spectrum---including repeated values---is
stronger than checking only the ground energy.  Degeneracies are part of the
expected answer, not noise to be discarded.

\section{Atomic limit and implementation agreement}

At $t=0$, every basis state is already an eigenstate and
\begin{equation}
E(s)=U\,\operatorname{popcount}(b_\uparrow\mathbin{\texttt{\&}}b_\downarrow).
\end{equation}
There must be no off-diagonal matrix elements.  Finally, for seeded random
vectors, require
\begin{equation}
\norm{H_{\mathrm{CSR}}x-H_{\mathrm{free}}x}_2
\end{equation}
to be consistent with floating-point roundoff, and require Eq.~\eqref{eq:residual}
for every computed eigenpair.

Run the complete executable validation with
\begin{verbatim}
.venv/bin/python -m pytest -q
\end{verbatim}
The tutorial programs supplement rather than replace these automated tests.

\paragraph{Checkpoint.}
A new representation should reproduce this entire ladder on the largest
shared system size before its performance or scaling claims are considered.

\chapter{General one-body hopping matrices}
\label{chap:generalhopping}

The nearest-neighbor chain is a specialization of
\begin{equation}
H(h)=\sum_{ij,\sigma}h_{ij}c_{i\sigma}^\dagger c_{j\sigma}
 +U\sum_i n_{i\uparrow}n_{i\downarrow}.
\label{eq:generalh}
\end{equation}
The interface convention is explicit: \code{h[i,j]} is the coefficient of
$c_i^\dagger c_j$.  The same $h$ is applied to both spin species.

\section{Validation and supported terms}

The input must have shape $(L,L)$, contain finite numerical values, and satisfy
\begin{equation}
h=h^\dagger
\end{equation}
to relative and absolute tolerance $10^{-12}$.  Accepted roundoff-level
asymmetry is projected away by storing $(h+h^\dagger)/2$.  The validated matrix
is copied and made read-only, preventing caller mutation from changing a
calculation after construction.

The entries can describe
\begin{itemize}
\item arbitrary-range or graph hopping;
\item real or complex Peierls phases;
\item spin-independent onsite potentials on the diagonal;
\item dense matrices produced by a one-particle basis transformation.
\end{itemize}
Each nonzero entry is evaluated through the same $c_i^\dagger c_j$ operator
path used earlier.  CSR and matrix-free complex-Hermitian actions are therefore
available without introducing a second sign convention.

For a ring threaded by total phase $\Phi$, one convenient gauge is
\begin{equation}
h_{i+1,i}=-t e^{i\Phi/L},\qquad
h_{i,i+1}=\left(h_{i+1,i}\right)^*.
\end{equation}

\lstinputlisting[caption={A complex flux model with longer-range hopping.},label={lst:generalh}]{code/step07_generic_hopping.py}

Complex Hermitian matrices produce real eigenvalues but generally complex
eigenvectors.  The generalized solver therefore preserves complex vector
dtype, uses a deterministic complex initial vector, and recomputes complex
residuals with the matrix-free action.

\section{The orbital-rotation caveat}

The general one-body interface is necessary for wavelets and orbital rotations,
but it is not sufficient for the interacting model.  Let a unitary matrix $R$
define new operators through
\begin{equation}
c_{i\sigma}=\sum_a R_{ia}d_{a\sigma}.
\end{equation}
The one-body matrix transforms simply:
\begin{equation}
h' = R^\dagger h R.
\end{equation}
The onsite interaction does not generally remain onsite.  Instead,
\begin{align}
U\sum_i n_{i\uparrow}n_{i\downarrow}
&=\sum_{abcd}V_{abcd}
d_{a\uparrow}^\dagger d_{b\uparrow}
d_{c\downarrow}^\dagger d_{d\downarrow},\\
V_{abcd}
&=U\sum_i R^*_{ia}R_{ib}R^*_{ic}R_{id}.
\label{eq:rotatedinteraction}
\end{align}
Therefore, passing only $h'$ to Eq.~\eqref{eq:generalh} while retaining a local
$U$ usually defines a different physical model.  It is exact at $U=0$ and in
special cases where the rotation preserves the local interaction.  A general
four-index two-body operator is the required next layer for faithful wavelet,
natural-orbital, or learned-orbital calculations at finite $U$.

\paragraph{Checkpoint.}
At $U=0$, compare the complete many-body spectrum obtained from a general $h$
with sums of its one-particle eigenvalues.  For complex $h$, additionally test
Hermiticity, CSR versus matrix-free action on complex random vectors, real
eigenvalues, and small residuals.

\chapter{From reference solver to research experiments}
\label{chap:extensions}

The current solver is intentionally small.  Its purpose is not to compete with
high-performance many-body packages, but to provide a transparent numerical
contract for new representations.

\section{A safe extension sequence}

The following order preserves independent validation at each stage.
\begin{enumerate}
\item \textbf{Arbitrary one-body graphs.}  Already supported through
      Chapter~\ref{chap:generalhopping}; use this for longer-range hopping, disorder,
      and flux.
\item \textbf{Natural-orbital diagnostics.}  Diagonalize the one-body density
      matrix from Eq.~\eqref{eq:obdm}.  At first, treat its eigenvectors as an
      analysis tool rather than silently changing the interacting Hamiltonian.
\item \textbf{General two-body tensors.}  Implement Eq.~\eqref{eq:rotatedinteraction}
      with an explicit index and operator-ordering convention.  Check unitary
      basis invariance of the complete small-system spectrum.
\item \textbf{Wavelet and learned rotations.}  Transform both one- and two-body
      terms.  Compare energies, residuals, reduced density matrices, and
      truncation errors against the site-basis result.
\item \textbf{Entanglement diagnostics.}  Add subsystem maps and reduced density
      matrices only after fixing the tensor-product ordering used for a cut.
\item \textbf{External methods.}  Compare DMRG, neural-network, and transformer
      wavefunctions on identical finite chains and boundary conventions.
\end{enumerate}

\section{A reproducible experiment contract}

Every comparison should record
\begin{itemize}
\item $L$, $(N_\uparrow,N_\downarrow)$, $t$, $U$, and boundary condition;
\item the spin-orbital and basis enumeration conventions;
\item Hilbert-space dimension and whether the action is CSR or matrix-free;
\item solver tolerance, random seed, eigenvalue count, and measured residuals;
\item all orbital transforms and whether the interaction was transformed;
\item observables with their normalization conventions;
\item agreement with at least one analytic limit and the site-basis reference.
\end{itemize}
Energy alone is not a sufficient comparison.  Two states with similar energy
can have different correlations, symmetry content, or reduced density
matrices.

\section{Performance interpretation}

The half-filled dimensions in Chapter~\ref{chap:basis} explain the practical
range.  On the reference host, $L=8$ is quick and $L=10$ is an upper practical
target; $L=12$ is generally too slow and memory-heavy in pure Python.  A dense
general hopping matrix also has up to $L^2$ terms per spin, so a matrix-free
product can be much slower than the nearest-neighbor action at the same
$\basisdim$.

Optimization should follow profiling and retain a slow trusted path.  Useful
future optimizations include cached transition tables, compiled bit kernels,
symmetry-sector reductions, and parallel observable evaluation.  Each optimized
path should continue to pass the random-vector and full-spectrum comparisons
against the current implementation.

\section{Suggested exercises}

\begin{enumerate}
\item Add a site-dependent diagonal potential through $h_{ii}$ and verify the
      $U=0$ full spectrum from the one-particle eigenvalues.
\item Add a next-nearest-neighbor edge whose source and destination have an
      occupied site between them; predict and test the sign.
\item Compute natural-orbital occupations for $U/t=0,4,16$ on the open $L=6$
      chain.  Check their sum and spin symmetry.
\item Construct a random unitary $R$ at $U=0$ and verify that $h$ and
      $R^\dagger hR$ give identical many-body spectra.
\item Implement the four-index interaction for $L=2$ and verify full spectral
      invariance under a random orbital rotation.
\end{enumerate}

The last exercise is the natural bridge from this tutorial to serious orbital
representation research: it tests not only code, but whether the transformed
Hamiltonian still represents the same physics.


\appendix
\chapter{Package map}

\begin{center}
\begin{tabular}{ll}
\toprule
Module & Responsibility \\
\midrule
\code{basis.py} & Fixed-sector enumeration and index maps \\
\code{operators.py} & Creation, annihilation, and hopping signs \\
\code{hamiltonian.py} & CSR, matrix-free, and generic one-body matrices \\
\code{lanczos.py} & Lowest eigenpairs and residual diagnostics \\
\code{observables.py} & Local, correlation, and Fourier observables \\
\code{analytic.py} & Two-site and noninteracting reference results \\
\bottomrule
\end{tabular}
\end{center}

The production source is intentionally compact.  Read each module after its
corresponding chapter, then inspect the matching tests to see how the stated
invariants are enforced.

\end{document}
